% !TEX root = ../presentation.tex
% =============================================================
%  Subsection — Series compensation of AC overhead lines
%  Part of ac-transmission; \input from ac-transmission.tex,
%  directly after natural-power.tex.
%
%  Figures in ./fig (sc_*.pdf), all vector, generated by
%  fig/series_compensation_figures.py from the exact two-port
%  equations with the capacitor bank cascaded as a series element
%  (numpy + matplotlib only).  Same model as the worked example in
%  the notebook ac_ohl_series_compensation.ipynb.
%
%  Reference line throughout, given in full on reference-line.tex:
%    500 kV, 60 Hz, 300 mi (483 km), flat,
%    3x Bluebird ACSR 2156 kcmil per phase, 18 in bundle, 35 ft phase
%    spacing, R at 75 C:
%    R' = 56.9 m$\Omega$/km, X' = 0.336 $\Omega$/km, C' = 12.7 nF/km,
%    X_L = 162 $\Omega$, R = 27.5 $\Omega$, Z_w = 266 $\Omega$, P_nat = 939 MW
% =============================================================
\providecommand{\acfig}[1]{ac-transmission/fig/#1}

\intermezzo{Series Compensation}{Shortening the electrical length of an OHL}

% ---------------------------------------------------------------
\begin{frame}{Series compensation --- principle}
  \setlength{\abovedisplayskip}{3pt}\setlength{\belowdisplayskip}{3pt}
  \noindent{\small A capacitor bank in series with the line cancels part of
  the series inductive reactance. The \emph{degree of compensation} $k$ is
  defined against the total series reactance of the line:}
  \[ k = \frac{X_C}{X_L}, \qquad X_L = \omega\cdot L'\cdot l,
     \qquad X_{eff} = \left(1-k\right)\cdot X_L \]
  \[ P = \frac{U^2}{\left(1-k\right)\cdot X_L}\cdot\sin\delta
     \quad\Rightarrow\quad
     P_{max} \approx \frac{U^2}{\left(1-k\right)\cdot X_L}
     \qquad\text{\small(lossless, } \left|U_S\right|=\left|U_R\right|=U\text{\small)} \]
  {\small
  \begin{itemize}
    \setlength{\itemsep}{0.2em}
    \item The line behaves as if it were shorter: for a given power the
          transmission angle $\delta$ falls and the steady-state stability
          margin rises.
    \item The rise $I\cdot X_C$ across the bank scales with the current:
          self-regulating, unlike shunt compensation.
    \item Practical range $k\approx0.25\ldots0.70$; above that SSR,
          distance-protection behaviour and bank cost dominate.
  \end{itemize}}
  \atnote{$X_C$ reactance of the bank at $\omega$; $l$ line length;
  $\delta$ angle between sending- and receiving-end voltages.}
\end{frame}

% ---------------------------------------------------------------
\begin{frame}{Series compensation --- single-line diagram}
  \centering
  \begin{tikzpicture}[x=1mm,y=1mm,scale=0.8,transform shape,
      wire/.style={draw=atInk,line width=0.35mm,line cap=round,line join=round},
      thin/.style={draw=atInk,line width=0.25mm,line cap=round,line join=round},
      bus/.style={draw=atInk,line width=0.9mm,line cap=butt},
      lbl/.style={font=\footnotesize,atInk},
      sub/.style={font=\scriptsize,atGray},
      dot/.style={fill=atInk,circle,inner sep=0pt,minimum size=1.1mm}]
    % buses and line
    \draw[bus] (0,-8) -- (0,8);   \node[lbl,above] at (0,8.5) {S};
    \draw[bus] (120,-8) -- (120,8); \node[lbl,above] at (120,8.5) {R};
    \draw[wire] (0,0) -- (30,0);
    \draw[wire] (90,0) -- (120,0);
    \node[sub] at (15,3.2) {$l/2$};
    \node[sub] at (105,3.2) {$l/2$};
    % line disconnectors into the platform
    \draw[wire] (30,0) -- (33,0);
    \draw[wire] (33,0) -- (38.5,3.5);            % DS open-style blade
    \draw[wire] (39,0) -- (42,0);
    \node[sub,above] at (36,4) {DS1};
    \draw[wire] (78,0) -- (81,0);
    \draw[wire] (81,0) -- (86.5,3.5);
    \draw[wire] (87,0) -- (90,0);
    \node[sub,above] at (84,4) {DS2};
    % main branch: series capacitor
    \node[dot] (a) at (42,0) {};  \node[dot] (b) at (78,0) {};
    \draw[wire] (42,0) -- (57,0);
    \draw[wire] (57,-4) -- (57,4);  \draw[wire] (60,-4) -- (60,4);
    \draw[wire] (60,0) -- (78,0);
    \node[lbl,above] at (58.5,4.5) {$C$\quad{\scriptsize\color{atGray}$X_C = k\cdot X_L$}};
    % MOV branch
    \draw[wire] (42,0) -- (42,-10) -- (54,-10);
    \draw[wire] (54,-13) rectangle (66,-7);
    \draw[thin] (55,-12) -- (65,-8);
    \draw[wire] (66,-10) -- (78,-10) -- (78,0);
    \node[sub,below] at (60,-13.3) {MOV};
    % damping reactor + bypass breaker branch
    \draw[wire] (42,-10) -- (42,-20) -- (50,-20);
    \draw[wire] (50,-22) rectangle (58,-18);     % IEC reactor (rectangle)
    \draw[wire] (58,-20) -- (63,-20);
    \draw[wire] (63,-22) -- (67,-18);  \draw[wire] (63,-18) -- (67,-22);   % CB as cross
    \draw[wire] (67,-20) -- (78,-20) -- (78,-10);
    \node[sub,below] at (54,-22.3) {$L_d$};
    \node[sub,below] at (65,-22.3) {bypass CB};
    % bypass disconnector around the whole bank
    \draw[wire] (33,0) -- (33,-33) -- (55,-33);
    \draw[wire] (55,-33) -- (60.5,-29.5);
    \draw[wire] (61,-33) -- (87,-33) -- (87,0);
    \node[sub,below] at (58,-34) {bypass DS};
    % platform outline
    \draw[atGray,densely dashed,thin] (39,-27.5) rectangle (81,11.5);
    \node[sub,anchor=south east] at (81,11.8) {capacitor platform, line potential};
    % current arrow
    \draw[-{Stealth[length=2mm]},atGray] (8,-3.5) -- (22,-3.5) node[midway,below,sub] {$I$};
  \end{tikzpicture}

  \vspace{0.2em}
  {\small
  \begin{itemize}
    \setlength{\itemsep}{0.15em}
    \item MOV limits the capacitor voltage under fault current; the bypass CB
          with damping reactor $L_d$ takes over at the MOV energy limit.
    \item DS1/DS2 isolate the platform; the bypass DS keeps the line in
          service with the bank out ($k=0$).
  \end{itemize}}
  \atnote{One phase shown; one platform per phase at the line midpoint of the
  reference case. Symbols IEC 60617.}
\end{frame}

% ---------------------------------------------------------------
\begin{frame}{Series compensation --- model and cases}
  \begin{columns}[T,onlytextwidth]
    \begin{column}{0.50\textwidth}
      {\small The two-port model of the natural-power slides is reused; the
      bank is cascaded at the midpoint as a series element
      $\begin{bmatrix}1 & -j\cdot X_C\\ 0 & 1\end{bmatrix}$ between the two
      half-lines.}\\[0.6em]
      {\small
      \begin{itemize}
        \setlength{\itemsep}{0.3em}
        \item $\beta\cdot l = 35.2^\circ$: at 300\,mi the line is well
              beyond the thermal plateau of the loadability curve.
        \item The compensation degree is swept over
              $k = 0,\;0.3,\;0.5,\;0.7$.
      \end{itemize}}
    \end{column}
    \begin{column}{0.47\textwidth}
      \renewcommand{\arraystretch}{1.15}
      {\footnotesize
      \begin{tabular}{@{}l l r@{}}
        \textbf{Quantity} & & \multicolumn{1}{c}{\textbf{Value}} \\
        \hline
        $X_L$ & $\Omega$ & 162.4 \\
        $R$ & $\Omega$ & 27.5 \\
        $Z_w$ & $\Omega$ & 266 \\
        $P_{nat}$ & MW & 939 \\
        \hline
      \end{tabular}}\\[0.4em]
      {\color{atGray}\footnotesize The line and its geometry are given on
      the reference-line slide; only the quantities the compensation uses
      are repeated here.}
    \end{column}
  \end{columns}
\end{frame}

% ---------------------------------------------------------------
\begin{frame}{Power--angle curves}
  \begin{columns}[c,onlytextwidth]
    \begin{column}{0.55\textwidth}\centering
      \topoimg[width=\linewidth,height=0.62\textheight,keepaspectratio]{\acfig{sc_power_angle.pdf}}
    \end{column}
    \begin{column}{0.43\textwidth}
      \renewcommand{\arraystretch}{1.15}
      {\footnotesize
      \begin{tabular}{@{}c r r r@{}}
        \textbf{$k$} & \multicolumn{1}{c}{\textbf{$X_C$}}
                     & \multicolumn{1}{c}{\textbf{$P_{max}$}}
                     & \multicolumn{1}{c}{\textbf{$\delta$ at 1500\,MW}} \\
                     & \multicolumn{1}{c}{$\Omega$} & \multicolumn{1}{c}{MW}
                     & \multicolumn{1}{c}{$^\circ$} \\
        \hline
        0   & 0     & 1363 & --- \\
        0.3 & 48.7  & 1744 & 50.1 \\
        0.5 & 81.2  & 2093 & 35.7 \\
        0.7 & 113.7 & 2407 & 25.3 \\
        \hline
      \end{tabular}}\\[0.45em]
      {\small
      \begin{itemize}
        \setlength{\itemsep}{0.3em}
        \item Uncompensated, 1500\,MW ($1.6\cdot P_{nat}$) cannot be carried
              at all; $k=0.5$ carries it at $36^\circ$.
        \item $P_{max}$ moves to smaller angles as $k$ grows --- the line gets
              electrically shorter.
      \end{itemize}}
    \end{column}
  \end{columns}
  \atnote{Both terminals at 1.0\,p.u.; exact two-port with losses and
  charging. The lossless $U^2/\!\left(\left(1-k\right)\cdot X_L\right)$
  would give 1540\,MW uncompensated and 3080\,MW at $k=0.5$ ---
  $R=27.5\,\Omega$ and the charging are not negligible at this length.}
\end{frame}

% ---------------------------------------------------------------
\begin{frame}{Voltage profile at heavy load}
  \begin{columns}[c,onlytextwidth]
    \begin{column}{0.55\textwidth}\centering
      \topoimg[width=\linewidth,height=0.62\textheight,keepaspectratio]{\acfig{sc_voltage_profile.pdf}}
    \end{column}
    \begin{column}{0.43\textwidth}
      {\small
      \begin{itemize}
        \setlength{\itemsep}{0.4em}
        \item 900\,MW at $\cos\varphi=0.95$ lagging delivered at the receiving
              end, $U_R$ held at 1.0\,p.u.
        \item Uncompensated the sending end must sit at 1.217\,p.u.\ ---
              far outside the operating band.
        \item The midpoint bank lifts the voltage by $I\cdot X_C$ in the
              direction of power flow: $U_S$ falls to 1.157, 1.127 and
              1.107\,p.u.\ for $k=0.3$, 0.5 and 0.7.
        \item The receiving half of the line is unchanged --- the bank acts
              only on what is upstream of it, which is why end-of-line
              banks and split banks are also used.
      \end{itemize}}
    \end{column}
  \end{columns}
\end{frame}

% ---------------------------------------------------------------
\begin{frame}{Bank rating and protection}
  \begin{columns}[T,onlytextwidth]
    \begin{column}{0.50\textwidth}
      \setlength{\abovedisplayskip}{4pt}\setlength{\belowdisplayskip}{4pt}
      {\small The bank is rated for the planned continuous line current, not
      for the conductor thermal limit:}
      \[ U_C = I_{rated}\cdot X_C, \qquad
         Q_C = 3\cdot I_{rated}^{2}\cdot X_C \]
      {\small
      \begin{itemize}
        \setlength{\itemsep}{0.3em}
        \item Fault-current voltage is limited by a metal-oxide varistor
              (MOV) across the bank, protective level $2.0\ldots2.5\cdot U_C$,
              with bypass breaker and damping circuit.
        \item Platform insulation at line potential; the bank carries the
              full fault current until the bypass closes.
        \item Emergency ratings (30\,min, 10\,s) are specified separately
              from the continuous rating.
      \end{itemize}}
    \end{column}
    \begin{column}{0.47\textwidth}
      \renewcommand{\arraystretch}{1.15}
      {\footnotesize
      \begin{tabular}{@{}l l r@{}}
        \textbf{Quantity} & & \multicolumn{1}{c}{\textbf{Value}} \\
        \hline
        $k$ & & 0.5 \\
        $X_C$ & $\Omega$ & 81.2 \\
        $C$ per phase & \textmu F & 32.7 \\
        $I_{rated}$ & A & 2000 \\
        \hline
        $U_C$ per phase & kV & 162 \\
        $Q_C$ three-phase & Mvar & 974 \\
        MOV level ($2.3\cdot U_C$) & kV & 373 \\
        \hline
      \end{tabular}}\\[0.4em]
      {\color{atGray}\footnotesize $I_{rated}$ chosen as $\approx1.15\times$
      the current at 1500\,MW, 500\,kV. Midpoint bank on the reference line.}
    \end{column}
  \end{columns}
\end{frame}

% ---------------------------------------------------------------
\begin{frame}{Sub-synchronous resonance --- screening}
  \begin{columns}[c,onlytextwidth]
    \begin{column}{0.55\textwidth}\centering
      \topoimg[width=\linewidth,height=0.52\textheight,keepaspectratio]{\acfig{sc_ssr_screening.pdf}}
    \end{column}
    \begin{column}{0.43\textwidth}
      \setlength{\abovedisplayskip}{4pt}\setlength{\belowdisplayskip}{4pt}
      {\small Series resonance of bank, line and source:}
      \[ f_{er} = f_0\cdot\sqrt{\frac{X_C}{X_L + X_S}} \]
      {\small
      \begin{itemize}
        \setlength{\itemsep}{0.3em}
        \item Rotor currents see the complement $f_0 - f_{er}$, which must
              stay clear of the turbine-generator torsional modes
              (10--45\,Hz for large steam units).
        \item $X_S = 25\,\Omega$ (15\,\% on 1500\,MVA): $k=0.5$ gives
              $f_{er}=39.5$\,Hz, complement 20.5\,Hz --- within 2\,Hz of
              the 22\,Hz example mode.
        \item A hit at screening level calls for a detailed SSR study or a
              different $k$; type-3 wind adds SSCI.
      \end{itemize}}
    \end{column}
  \end{columns}
  \atnote{Torsional modes shown are illustrative; real values come from the
  generator manufacturer.}
\end{frame}
